% Seed 0 generous-parser recovery by extraction rule. Requires no LaTeX packages.
% Counts from analysis/generous-seed-0/scores.jsonl: rows that became valid under
% the permissive parser but were invalid under the original scorer.
\begin{table}[htbp]
\centering
\footnotesize
\setlength{\tabcolsep}{6pt}
\renewcommand{\arraystretch}{1.3}
\caption{Additional submissions recovered by each permissive extraction rule (seed~0).
Each row is one rule of the generous parser applied to submissions the original scorer
rejected. Valid means a readable submission was identified; correct means it passed the
benchmark's scoring. Rules operate only on previously invalid responses.}
\label{tab:generous-parser-recovery}
\begin{tabular}{@{}p{0.28\textwidth}rrp{0.42\textwidth}@{}}
\hline
Extraction rule & \multicolumn{1}{r}{Valid} & \multicolumn{1}{r}{Correct} & What the rule accepts \\
\hline
Extract code block & 221 & 33 & The last fenced or plain code block containing the
required signature, even when prose, Markdown, or reasoning surrounds it. The strict
scorer only read the final fenced block of the post-thinking text. \\
Insert definition keyword & 164 & 16 & A function body whose header omitted the
definition keyword (e.g.\ a bare \texttt{name(args):} line) but otherwise matched the
required entry point; the keyword is restored before parsing. \\
Leading option letter & 69 & 26 & A GPQA response that opens with a single choice
letter, optionally parenthesized, when no competing option label or reconsideration
language follows. \\
Marked answer line & 16 & 5 & Answer statements in looser forms than the strict
\texttt{ANSWER:} pattern, such as ``the answer is (B)'' or ``correct choice: C'',
including when they appear before the final line. \\
Leading integer & 6 & 0 & An AIME response whose first line is a standalone integer,
when the text shows no reconsideration. \\
Final standalone answer & 2 & 0 & A response whose entire final line is a bare answer
token, without any marker. \\
\hline
\textbf{Total} & \textbf{478} & \textbf{80} & \\
\hline
\end{tabular}
\end{table}
